% chapters/04/praktikum.tex
\begin{praktikumbox}
Praktikum ini mengikuti Tugas Praktikum 02 pada RPS: membuat siklus permainan
yang lengkap, membatasinya pada 60 bingkai per detik, menampilkan laju bingkai
pada judul jendela, lalu mengubah batas tersebut dan mencatat pengaruhnya
terhadap beban prosesor. Isilah lembar kerja pada
Tabel~\ref{tab:lembar-praktikum-02}.

\begin{langkah}
  \item \textbf{Jalankan siklus lengkap.} Ketik ulang
        \code{ch04_game_loop.py} dari Daftar~\ref{lst:ch04-game-loop} ke dalam
        folder \code{code}. Jalankan program, lalu amati tiga hal: judul
        jendela, batang penanda di bagian bawah, dan nilai laju terukur.

  \item \textbf{Ubah batas laju dari dalam program.} Tekan tombol panah atas
        dan bawah beberapa kali. Catat laju terukur pada setiap nilai batas.
        Judul jendela berubah seketika.

  \item \textbf{Bandingkan batas laju dengan pengukuran.} Jalankan
        \code{ch04_ukur_laju.py} dari Daftar~\ref{lst:ch04-ukur-laju} tanpa
        argumen. Catat seluruh baris tabel yang dihasilkan.

  \item \textbf{Amati pengaruhnya pada prosesor.} Buka pengelola tugas pada
        sistem operasi Anda (misalnya \istilah{Task Manager} pada Windows).
        Jalankan program siklus dengan batas 60, lalu dengan pilihan
        \code{--fps 0} yang berarti tanpa pembatasan. Amati pemakaian prosesor
        pada kedua keadaan dan catat perbedaannya.

  \item \textbf{Tambahkan satu penanda keadaan.} Tambahkan satu baris teks pada
        jendela yang menampilkan jumlah bingkai yang telah dijalankan sejak
        penekanan spasi terakhir. Uji apakah penanda itu benar-benar kembali ke
        nol ketika spasi ditekan.

  \item \textbf{Tulis catatan praktikum} pada folder \code{catatan} yang
        memuat hasil pengukuran Anda dan kesimpulan mengenai batas laju yang
        Anda pilih.

  \item \textbf{Simpan pekerjaan} ke repositori Git dengan pesan penyimpanan
        yang menerangkan isinya.
\end{langkah}
\end{praktikumbox}

\begin{table}[htbp]
\centering
\caption{Lembar kerja Praktikum 02}
\label{tab:lembar-praktikum-02}
{\setlength{\tabcolsep}{4pt}\footnotesize%
\begin{tabularx}{\textwidth}{@{}C{1.6cm} L{5.0cm} Y@{}}
\toprule
\headrow \textbf{Langkah} & \textbf{Hasil yang dicatat} & \textbf{Catatan Anda} \\
\midrule
1 & Laju terukur pada batas laju bawaan & \\[1.2em]
\addlinespace
2 & Laju terukur pada tiga nilai batas yang berbeda & \\[2.4em]
\addlinespace
3 & Baris tabel hasil pengukuran beserta pemakaian prosesornya & \\[2.4em]
\addlinespace
4 & Perbandingan pemakaian prosesor pada batas 60 dan tanpa batas & \\[2.4em]
\addlinespace
5 & Jumlah bingkai yang tercatat sebelum dan sesudah spasi ditekan & \\[1.2em]
\addlinespace
6 & Batas laju yang Anda pilih beserta alasannya & \\[2.0em]
\bottomrule
\end{tabularx}}
\end{table}

\begin{catatabox}
Pemakaian prosesor yang tampak pada pengelola tugas akan berbeda dari angka
pada Tabel~\ref{tab:hasil-pengukuran-laju}, karena pengukuran di buku
dilakukan tanpa jendela sedangkan program Anda menampilkan jendela
sungguhan. Yang perlu Anda bandingkan adalah perbandingan antara keadaan
berbatas dan tanpa batas, bukan angka mutlaknya.
\end{catatabox}
